Fix \(P\in\mathcal P_L\) and \(t\in\{0,1\}\), and suppose that \(a,a'\in[0,1]\) both satisfy the displayed point-anchor envelope.  This proof uses the shrinking-support argument.  For every \(H\in(0,1/2]\), assumption \(\mathrm{ass:uniform\text{-}disc\text{-}design}\) gives
\[
 P_X\{x\in\mathcal X_{t,P}:\lVert x-b\rVert_2<H\}
 =\mathbb P_P(T=t,R<H)\ge c_0H^2>0.
\]
The two envelope inequalities fail only on the union of two \(P_X\)-null sets.  Hence the positive-probability set above contains a point \(x_H\) at which both hold.  Assumption \(\mathrm{ass:radial\text{-}lipschitz}\) and the triangle inequality then yield
\[
 |a-a'|
 \le |a-\mu_t(x_H)|+|\mu_t(x_H)-a'|
 \le 2L\lVert x_H-b\rVert_2
 <2LH.
\]
Because this inequality holds for every \(H\in(0,1/2]\), letting \(H\downarrow0\) gives \(|a-a'|=0\), hence \(a=a'\).  The member properties in \(\mathrm{def:half\text{-}disc\text{-}class}\) assert existence of the declared anchor \(\theta_t\); the argument proves its uniqueness.  Therefore \(\theta_1-\theta_0\), and thus \(\tau_P(b)\), is a well-defined functional of the full-data law rather than a choice of conditional-mean version.